\newpage
\chapter{TINJAUAN PUSTAKA} \label{Bab II}

\section{Tinjauan Pustaka} \label{II.Tinjauan}
Penelitian tentang klasifikasi citra media sosial untuk respons bencana
telah berkembang pesat dalam beberapa tahun terakhir, khususnya sejak
tersedianya \textit{dataset} standar seperti \textit{CrisisMMD} yang
memungkinkan perbandingan yang adil antar berbagai metode. \par

\begin{enumerate}[noitemsep]
    \item Ofli et al. mengeksplorasi pendekatan multimodal yang
    menggabungkan informasi visual dan tekstual untuk meningkatkan
    akurasi klasifikasi pada \textit{dataset} \textit{CrisisMMD}
    \cite{ofli2020analysis}. Penelitian ini mengevaluasi tiga
    konfigurasi pelatihan: unimodal teks, unimodal citra, dan
    multimodal (teks dan citra). Untuk Task 1 (\textit{Informative
    Classification}), pendekatan unimodal berbasis citra mencapai
    akurasi 83.3\%, sedangkan kombinasi teks dan gambar meningkatkan
    akurasi menjadi 84.4\%. Untuk Task 2 (\textit{Humanitarian
    Categories}), unimodal citra mencapai akurasi 76.8\%, sementara
    pendekatan multimodal menghasilkan akurasi 78.4\%. Meskipun
    pendekatan multimodal memberikan peningkatan performa, angka
    unimodal berbasis citra tersebut menjadi \textit{baseline} yang
    relevan bagi penelitian yang berfokus pada modalitas citra saja.

    \item Mohanty et al. membandingkan beberapa arsitektur CNN
    (\textit{Inception V3}, VGG, dan ResNet) untuk klasifikasi citra
    badai secara biner (\textit{hurricane} vs.\ \textit{non-hurricane})
    \cite{mohanty2021multimodal}. Penelitian ini menunjukkan bahwa
    \textit{Tuned Inception V3} hasil \textit{fine-tuning} berbasis
    \textit{transfer learning} mencapai \textit{F1-score} 0.95 dan
    unggul sekitar 6--7\% dibandingkan arsitektur lainnya. Temuan ini
    mengonfirmasi pentingnya \textit{fine-tuning} dan \textit{transfer
    learning} dalam meningkatkan performa model CNN pada domain
    klasifikasi citra bencana, meskipun penelitian tersebut tidak
    mengeksplorasi pendekatan \textit{ensemble} maupun analisis
    efisiensi komputasi.

    \item Firmansyah et al. melakukan evaluasi sistematis terhadap lima
    arsitektur CNN (\textit{VGG16}, \textit{DenseNet201},
    \textit{MobileNetV2}, \textit{InceptionV3}, dan \textit{ResNet50})
    pada \textit{dataset} \textit{CrisisMMD}
    \cite{firmansyah2022ensemble}. Penelitian ini menemukan bahwa
    \textit{ensemble} sederhana menggunakan \textit{simple averaging}
    dari dua model terbaik (DenseNet201 dan VGG16, atau DenseNet201 dan
    MobileNetV2) memberikan peningkatan akurasi sekitar 1\%
    dibandingkan model tunggal terbaik, dengan akurasi mencapai 84.6\%.
    Namun, penelitian tersebut hanya menggunakan \textit{homogeneous
    ensemble} dari CNN dan tidak menganalisis pertukaran antara
    peningkatan akurasi dengan biaya komputasi seperti waktu inferensi
    dan penggunaan memori.

    \item Long et al. memperkenalkan \textit{CrisisViT}, sebuah
    arsitektur \textit{Vision Transformer} yang dirancang khusus untuk
    klasifikasi citra krisis \cite{long2023crisisvit}. Penelitian ini
    menunjukkan bahwa arsitektur berbasis \textit{Transformer} mulai
    diterapkan pada domain respons bencana dan dapat mencapai performa
    yang kompetitif. Namun, penelitian tersebut tidak melakukan
    komparasi langsung dengan arsitektur CNN modern dan tidak
    menganalisis pertukaran antara akurasi dengan efisiensi komputasi.

    \item Debnath et al. memperkenalkan \textit{ConvSFNet}, sebuah
    arsitektur berbasis \textit{ConvNeXt} yang dikombinasikan dengan
    blok \textit{Squeeze-and-Excitation} dan struktur \textit{Feature
    Pyramid Network} untuk klasifikasi citra bencana pada dataset MEDIC
    \cite{debnath2024convsfnet}. Penelitian ini membandingkan
    \textit{ConvSFNet} terhadap beberapa arsitektur dasar termasuk
    \textit{DenseNet121}, \textit{EfficientNet-B1}, \textit{ResNet50},
    dan \textit{VGG16} pada empat tugas klasifikasi yang mencakup
    \textit{Humanitarian} dan \textit{Informative}. Hasilnya
    menunjukkan bahwa \textit{ConvSFNet} meningkatkan rata-rata
    \textit{F1-score} sebesar 2.41\% dibandingkan seluruh model
    pembanding. Temuan ini mengonfirmasi bahwa arsitektur berbasis
    \textit{ConvNeXt} memiliki potensi yang kuat untuk klasifikasi
    citra bencana, meskipun penelitian tersebut tidak mengeksplorasi
    pendekatan \textit{ensemble} maupun arsitektur \textit{Transformer}.

    \item Ganaie et al. menyajikan tinjauan komprehensif terhadap
    berbagai pendekatan \textit{ensemble deep learning}, mencakup
    kategori \textit{bagging}, \textit{boosting}, \textit{stacking},
    \textit{homogeneous ensemble}, \textit{heterogeneous ensemble},
    serta strategi \textit{decision fusion} \cite{ganaie2022ensemble}.
    Penelitian ini secara khusus membahas bahwa \textit{heterogeneous
    ensemble} yang menggabungkan arsitektur yang berbeda secara
    mendasar cenderung menghasilkan keragaman prediksi yang lebih
    tinggi, yang pada akhirnya berkontribusi pada peningkatan performa
    dibandingkan \textit{homogeneous ensemble}. Temuan ini memberikan
    landasan teoritis yang kuat bagi penggunaan \textit{ensemble}
    heterogen yang menggabungkan CNN dan \textit{Transformer} dalam
    penelitian ini.
\end{enumerate}

Berdasarkan tinjauan terhadap penelitian-penelitian terdahulu, terdapat
beberapa celah yang belum terisi. Pertama, studi komparatif yang
mencakup ragam arsitektur dari CNN murni hingga \textit{Transformer}
murni secara sistematis pada \textit{dataset CrisisMMD} belum pernah
dilakukan. Kedua, potensi \textit{heterogeneous ensemble} yang
menggabungkan CNN dan \textit{Transformer} untuk domain klasifikasi
informasi krisis belum dieksplorasi, padahal secara teoritis keragaman
arsitektur yang lebih tinggi berpotensi menghasilkan peningkatan
performa yang lebih signifikan. Ketiga, analisis sistematis terhadap
faktor-faktor yang mempengaruhi performa \textit{ensemble} dalam bentuk
\textit{ablation study} masih sangat jarang dilakukan pada domain ini.
Penelitian ini mengisi ketiga celah tersebut dengan melakukan analisis
komparatif sistematis antara empat arsitektur (\textit{EfficientNetV2-M},
\textit{ConvNeXt-Base}, \textit{Swin-Small}, dan ViT-B/16), berbagai
konfigurasi \textit{ensemble} heterogen, serta \textit{ablation study}
pada ketiga tugas klasifikasi \textit{CrisisMMD}. \par

\begin{longtable}{|b{0.04\textwidth}|p{0.22\textwidth}|p{0.18\textwidth}
|p{0.18\textwidth}|p{0.22\textwidth}|}
    \caption{Ringkasan Penelitian Terdahulu}
    \label{table:2.literasi}\\
    \hline
    \textbf{No.} & \textbf{Judul} & \textbf{Masalah} &
    \textbf{Metode} & \textbf{Hasil} \\
    \hline
    \endfirsthead
    \hline
    \textbf{No.} & \textbf{Judul} & \textbf{Masalah} &
    \textbf{Metode} & \textbf{Hasil} \\
    \hline
    \endhead
    1. &
    Analysis of social media data using multimodal deep learning for
    disaster response (Ofli et al., 2020) \cite{ofli2020analysis} &
    Klasifikasi otomatis citra bencana dari media sosial dengan
    memanfaatkan modalitas teks dan citra secara bersamaan &
    CNN + teks multimodal (unimodal teks, unimodal citra, multimodal)
    pada \textit{CrisisMMD} &
    T1: 84.4\% (multimodal), 83.3\% (image-only); T2: 78.4\%
    (multimodal), 76.8\% (image-only). Tidak ada analisis
    \textit{trade-off} komputasi \\
    \hline
    2. &
    A multi-modal approach towards mining social media data during
    natural disasters (Mohanty et al., 2021) \cite{mohanty2021multimodal} &
    Klasifikasi citra badai biner (\textit{hurricane} vs.\
    \textit{non-hurricane}) tanpa eksplorasi \textit{ensemble} &
    Perbandingan CNN (\textit{Inception V3}, VGG, ResNet) dengan
    \textit{transfer learning fine-tuning} pada dataset Hurricane Irma &
    \textit{Tuned Inception V3} mencapai F1-score 0.95, unggul 6--7\%
    dari arsitektur lain. Tidak ada analisis efisiensi komputasi \\
    \hline
    3. &
    Ensemble learning for the classification of social media data in
    disaster response (Firmansyah et al., 2022)
    \cite{firmansyah2022ensemble} &
    Peningkatan akurasi klasifikasi citra bencana pada
    \textit{CrisisMMD} menggunakan \textit{homogeneous ensemble} CNN &
    \textit{Ensemble} lima CNN (VGG16, DenseNet201, MobileNetV2,
    InceptionV3, ResNet50) dengan \textit{simple averaging} &
    Akurasi 84.6\% (DenseNet201 + VGG16). Tidak ada arsitektur
    \textit{Transformer} dan tidak ada analisis biaya komputasi \\
    \hline
    4. &
    CrisisViT: A robust vision transformer for crisis image
    classification (Long et al., 2023) \cite{long2023crisisvit} &
    Penerapan arsitektur \textit{Transformer} pada klasifikasi citra
    krisis tanpa komparasi langsung dengan CNN modern &
    \textit{CrisisViT} berbasis \textit{Vision Transformer} pada
    \textit{dataset} krisis &
    Performa kompetitif pada klasifikasi krisis. Tidak ada komparasi
    CNN dan tidak ada analisis \textit{trade-off} \\
    \hline
    5. &
    ConvSFNet: A novel architecture for disaster image analysis
    (Debnath et al., 2024) \cite{debnath2024convsfnet} &
    Klasifikasi citra bencana pada dataset MEDIC menggunakan arsitektur
    CNN modern tanpa eksplorasi \textit{ensemble} &
    \textit{ConvSFNet} berbasis \textit{ConvNeXt} + SE block + FPN,
    dibandingkan DenseNet121, EfficientNet-B1, ResNet50 &
    Rata-rata F1-score meningkat 2.41\% dibandingkan semua
    \textit{baseline}. Tidak ada \textit{Transformer} dan tidak ada
    \textit{ensemble} \\
    \hline
    6. &
    Ensemble deep learning: A review (Ganaie et al., 2022)
    \cite{ganaie2022ensemble} &
    Kurangnya tinjauan komprehensif tentang strategi
    \textit{ensemble deep learning} termasuk \textit{heterogeneous
    ensemble} &
    Tinjauan sistematis mencakup \textit{bagging}, \textit{boosting},
    \textit{stacking}, \textit{homogeneous} dan \textit{heterogeneous
    ensemble} &
    \textit{Heterogeneous ensemble} terbukti menghasilkan keragaman
    prediksi lebih tinggi dan performa lebih baik dibandingkan
    \textit{homogeneous ensemble} \\
    \hline
\end{longtable}


% ============================================================
\section{Dasar Teori} \label{II.Teori}

\subsection{Respons Bencana dan Informasi Krisis}
\label{II.Teori.Bencana}

\subsubsection{Siklus Manajemen Bencana}
Manajemen bencana umumnya terdiri dari empat fase utama: mitigasi,
kesiapsiagaan (\textit{preparedness}), respons, dan pemulihan
(\textit{recovery}). Fase respons adalah periode kritis dimana tindakan
segera diperlukan untuk menyelamatkan nyawa dan mengurangi dampak
bencana. Dalam fase ini, informasi yang akurat dan tepat waktu tentang
lokasi korban, tingkat kerusakan, dan kebutuhan bantuan menjadi sangat
penting untuk pengambilan keputusan yang efektif oleh tim respons
darurat. \par

\subsubsection{Peran Media Sosial dalam Respons Bencana}
Media sosial telah menjadi sumber informasi penting dalam respons
bencana karena kemampuannya menyediakan data hampir \textit{real-time}
dari masyarakat terdampak. Platform seperti \textit{Twitter} dan
\textit{Facebook} memungkinkan masyarakat untuk membagikan foto, video,
dan informasi tentang situasi di lapangan dalam hitungan menit setelah
kejadian. Informasi ini dapat memberikan \textit{situational awareness}
kepada tim respons, membantu identifikasi area yang membutuhkan bantuan
segera, dan memfasilitasi koordinasi upaya penyelamatan. Namun, volume
data yang sangat besar dan adanya informasi yang tidak relevan membuat
kurasi manual menjadi tidak praktis, sehingga diperlukan sistem
klasifikasi otomatis. \par

\subsubsection{Klasifikasi Informasi Krisis}
Klasifikasi informasi krisis bertujuan untuk memisahkan informasi yang
berguna untuk respons bencana dari informasi yang tidak relevan. Tugas
klasifikasi utama meliputi: (1) identifikasi apakah suatu citra bersifat
informatif atau tidak untuk respons bencana, (2) kategorisasi jenis
informasi kemanusiaan yang terkandung seperti kerusakan infrastruktur,
individu terdampak, dan operasi penyelamatan, serta (3) penilaian
tingkat kerusakan yang ditampilkan dalam citra. Klasifikasi otomatis
yang akurat dapat secara signifikan mempercepat proses triase informasi
dan membantu tim respons fokus pada informasi yang paling kritis. \par


\subsection{Deep Learning} \label{II.Teori.DL}
\textit{Deep learning} adalah cabang dari \textit{machine learning}
yang menggunakan arsitektur jaringan syaraf tiruan berlapis banyak untuk
mempelajari representasi data secara hierarkis dari data mentah tanpa
memerlukan rekayasa fitur manual \cite{lecun2015deep}. Setiap lapisan
dalam jaringan mempelajari representasi yang semakin abstrak dari data
masukan, sehingga model dapat menangkap pola yang kompleks dan beragam.
Keberhasilan \textit{deep learning} didorong oleh tiga faktor utama:
ketersediaan data berlabel dalam jumlah besar, peningkatan kapasitas
komputasi melalui GPU (\textit{Graphics Processing Unit}), dan
pengembangan teknik pelatihan yang lebih efektif
\cite{lecun2015deep}. \par

Dalam konteks \textit{computer vision}, \textit{deep learning} telah
merevolusi cara mesin memahami dan menginterpretasi citra. Klasifikasi
citra adalah tugas fundamental yang bertujuan untuk mengkategorikan
citra ke dalam kelas-kelas yang telah ditentukan berdasarkan konten
visualnya \cite{lecun1998gradient}. Dalam konteks respons bencana,
klasifikasi citra bertujuan untuk mengidentifikasi informasi yang
relevan dari ribuan gambar yang dibagikan di media sosial selama
kejadian bencana, mencakup pembedaan citra informatif dari yang tidak
informatif, kategorisasi jenis informasi kemanusiaan, serta penilaian
tingkat kerusakan. \par


\subsection{Convolutional Neural Network} \label{II.Teori.CNN}

\subsubsection{Arsitektur CNN}
\textit{Convolutional Neural Network} (CNN) adalah arsitektur
\textit{deep learning} yang dirancang khusus untuk memproses data
berbentuk grid seperti citra \cite{lecun1998gradient}. CNN terdiri dari
beberapa komponen utama: \textit{convolutional layers} yang mengekstraksi
fitur lokal melalui operasi konvolusi dengan bobot yang dapat dipelajari,
\textit{pooling layers} yang mengurangi dimensi spasial sambil
mempertahankan informasi penting, dan \textit{fully connected layers}
yang melakukan klasifikasi berdasarkan fitur yang telah diekstraksi. \par

Keunggulan utama CNN terletak pada \textit{inductive bias} (asumsi
bawaan arsitektur) berupa lokalitas dan invarians translasi yang sesuai
dengan karakteristik data citra. \textit{Layer} awal biasanya
mendeteksi fitur-fitur sederhana seperti tepi dan tekstur, sementara
\textit{layer} yang lebih dalam menggabungkan fitur-fitur tersebut
menjadi representasi yang lebih kompleks dan abstrak. Arsitektur ini
telah menjadi standar untuk berbagai tugas \textit{computer vision}
sejak keberhasilan \textit{AlexNet} pada kompetisi \textit{ImageNet}
tahun 2012 \cite{krizhevsky2012alexnet}. Dalam domain klasifikasi
citra untuk respons bencana, Nguyen et al. merupakan penelitian pionir
yang mendemonstrasikan bahwa arsitektur CNN seperti \textit{AlexNet}
dan VGG mampu mengidentifikasi informasi kritis terkait bencana dari
media sosial dengan akurasi yang tinggi \cite{nguyen2017robust}. \par

% TODO: Tambahkan gambar arsitektur CNN (conv → pool → FC)
% Sumber: LeCun et al. (1998) atau Krizhevsky et al. (2012)
\begin{figure}[h]
    \centering
    % \includegraphics[width=0.85\textwidth]{gambar/cnn_architecture.png}
    \fbox{\rule{0pt}{4cm}\rule{0.85\textwidth}{0pt}}
    \caption{Arsitektur umum \textit{Convolutional Neural Network}
    terdiri dari lapisan konvolusi, \textit{pooling}, dan
    \textit{fully connected} \cite{lecun1998gradient}}
    \label{fig:cnn_architecture}
\end{figure}

\subsubsection{EfficientNetV2}
\textit{EfficientNet} diperkenalkan oleh Tan dan Le sebagai keluarga
arsitektur CNN yang berfokus pada efisiensi melalui \textit{compound
scaling} \cite{tan2019efficientnet}. Berbeda dengan pendekatan
tradisional yang hanya meningkatkan kedalaman atau lebar jaringan,
\textit{EfficientNet} menggunakan \textit{compound coefficient} untuk
menyeimbangkan peningkatan kedalaman, lebar, dan resolusi input secara
bersamaan, menghasilkan model yang lebih akurat dengan jumlah parameter
yang lebih sedikit. \par

\textit{EfficientNetV2} merupakan perkembangan lanjutan yang
diperkenalkan oleh Tan dan Le dengan fokus pada peningkatan kecepatan
pelatihan dan akurasi \cite{tan2021efficientnetv2}. Perbedaan utama
dari generasi sebelumnya adalah penggunaan blok \textit{Fused-MBConv}
pada lapisan awal jaringan yang menggabungkan operasi konvolusi biasa
dan \textit{depthwise convolution} menjadi satu operasi konvolusi tunggal,
menghasilkan kecepatan pelatihan yang lebih tinggi. Selain itu,
\textit{EfficientNetV2} menggunakan strategi \textit{progressive
learning} yang secara bertahap meningkatkan ukuran citra masukan dan
kekuatan augmentasi selama pelatihan. Varian \textit{EfficientNetV2-M} yang digunakan dalam penelitian ini
memiliki sekitar 54 juta parameter dengan FLOPs sekitar 24 GFLOPs,
menjadikannya representasi CNN modern yang efisien dalam spektrum
arsitektur yang dibandingkan \cite{tan2021efficientnetv2}. \par

% TODO: Tambahkan gambar perbandingan blok MBConv vs Fused-MBConv
% Sumber: Tan & Le (2021), Figure 2
\begin{figure}[h]
    \centering
    % \includegraphics[width=0.75\textwidth]{gambar/efficientnetv2_block.png}
    \fbox{\rule{0pt}{3.5cm}\rule{0.75\textwidth}{0pt}}
    \caption{Perbandingan blok \textit{MBConv} (EfficientNet generasi
    pertama) dan \textit{Fused-MBConv} (EfficientNetV2)
    \cite{tan2021efficientnetv2}}
    \label{fig:efficientnetv2_block}
\end{figure}

\subsubsection{ConvNeXt}
\textit{ConvNeXt} diperkenalkan oleh Liu et al. sebagai arsitektur CNN
yang dirancang ulang dengan mengadopsi prinsip-prinsip desain dari
arsitektur \textit{Transformer} ke dalam kerangka konvolusional
\cite{liu2022convnext}. Motivasi utama penelitian ini adalah untuk
menjawab pertanyaan sejauh mana CNN murni dapat bersaing dengan
\textit{Vision Transformer} jika menerapkan praktik desain yang modern.
Beberapa pembaruan kunci yang diadopsi mencakup penggunaan
\textit{depthwise convolution} berukuran besar (7$\times$7) yang
memungkinkan setiap filter menangkap konteks spasial yang lebih luas,
lapisan normalisasi \textit{LayerNorm} (normalisasi per lapisan)
sebagai pengganti \textit{BatchNorm} (normalisasi per batch), fungsi
aktivasi \textit{GELU} sebagai pengganti \textit{ReLU}, serta
pengurangan jumlah lapisan aktivasi dan normalisasi per blok. \par

Hasil evaluasi menunjukkan bahwa \textit{ConvNeXt} mampu melampaui
\textit{Swin Transformer} pada berbagai tolok ukur \textit{computer
vision} dengan kompleksitas komputasi yang sebanding
\cite{liu2022convnext}. Hal ini menunjukkan bahwa batas kemampuan CNN
belum sepenuhnya tercapai dan prinsip konvolusional masih relevan
ketika dikombinasikan dengan desain arsitektur yang modern. Varian
\textit{ConvNeXt-Base} yang digunakan dalam penelitian ini memiliki
sekitar 88 juta parameter dengan FLOPs sekitar 15 GFLOPs
\cite{liu2022convnext}. \par

% TODO: Tambahkan gambar blok desain ConvNeXt vs ResNet
% Sumber: Liu et al. (2022), Figure 3
\begin{figure}[h]
    \centering
    % \includegraphics[width=0.75\textwidth]{gambar/convnext_block.png}
    \fbox{\rule{0pt}{3.5cm}\rule{0.75\textwidth}{0pt}}
    \caption{Evolusi desain blok dari \textit{ResNet} menuju
    \textit{ConvNeXt} melalui modernisasi bertahap
    \cite{liu2022convnext}}
    \label{fig:convnext_block}
\end{figure}


\subsection{Vision Transformer} \label{II.Teori.ViT}

\subsubsection{Arsitektur Transformer}
Arsitektur \textit{Transformer} awalnya diperkenalkan untuk pemrosesan
bahasa natural oleh Vaswani et al. dengan mekanisme inti berupa
\textit{multi-head self-attention} (atensi mandiri multi-kepala)
\cite{vaswani2017attention}. Berbeda dengan CNN yang memproses
informasi secara lokal melalui operasi konvolusi, mekanisme
\textit{self-attention} memungkinkan setiap elemen dalam urutan masukan
untuk berinteraksi secara langsung dengan semua elemen lainnya,
sehingga mampu menangkap ketergantungan jarak jauh dalam satu langkah.
Setiap kepala atensi (\textit{attention head}) mempelajari hubungan
yang berbeda antar elemen, dan hasilnya digabungkan untuk menghasilkan
representasi yang kaya. Komponen utama lainnya adalah
\textit{position-wise feed-forward network} yang diterapkan secara
independen pada setiap posisi, serta \textit{residual connections}
(koneksi sisa) dan \textit{layer normalization} yang memfasilitasi
pelatihan jaringan yang dalam \cite{vaswani2017attention}. \par

% TODO: Tambahkan gambar arsitektur Transformer (multi-head self-attention)
% Sumber: Vaswani et al. (2017), Figure 2
\begin{figure}[h]
    \centering
    % \includegraphics[width=0.55\textwidth]{gambar/transformer_attention.png}
    \fbox{\rule{0pt}{5cm}\rule{0.55\textwidth}{0pt}}
    \caption{Mekanisme \textit{multi-head self-attention} pada arsitektur
    \textit{Transformer} \cite{vaswani2017attention}}
    \label{fig:transformer_attention}
\end{figure}

\subsubsection{Vision Transformer (ViT-B/16)}
\textit{Vision Transformer} (ViT) adalah adaptasi dari arsitektur
\textit{Transformer} ke domain \textit{computer vision}
\cite{dosovitskiy2021image}. ViT membagi citra menjadi
\textit{patch-patch} berukuran tetap (16$\times$16 piksel untuk
ViT-B/16), kemudian mem-\textit{flatten} setiap \textit{patch} menjadi
vektor dan memprosesnya sebagai urutan token menggunakan mekanisme
\textit{self-attention}. Sebuah token \textit{class} yang dapat
dipelajari ditambahkan pada awal urutan dan digunakan sebagai representasi
keseluruhan citra untuk klasifikasi akhir. \par

Mekanisme \textit{self-attention} pada ViT memungkinkan setiap
\textit{patch} untuk berinteraksi dengan semua \textit{patch} lainnya
sejak \textit{layer} pertama, memberikan kemampuan untuk menangkap
ketergantungan jarak jauh tanpa dibatasi oleh \textit{receptive field}
(area citra yang dapat dilihat oleh satu neuron) seperti pada CNN
\cite{dosovitskiy2021image}. \textit{Positional encoding} (penyandian
posisi) ditambahkan untuk mempertahankan informasi posisi spasial dari
setiap \textit{patch}. ViT membutuhkan data \textit{pre-training} yang
sangat besar untuk mencapai performa optimal karena tidak memiliki
\textit{inductive bias} bawaan seperti CNN. ViT-B/16 yang digunakan
dalam penelitian ini memiliki 12 \textit{layer Transformer encoder}
dengan dimensi tersembunyi 768 dan 12 \textit{attention heads},
menghasilkan total sekitar 86 juta parameter
\cite{dosovitskiy2021image}. \par

% TODO: Tambahkan gambar arsitektur ViT (patch embedding + token sequence)
% Sumber: Dosovitskiy et al. (2021), Figure 1
\begin{figure}[h]
    \centering
    % \includegraphics[width=0.85\textwidth]{gambar/vit_architecture.png}
    \fbox{\rule{0pt}{4cm}\rule{0.85\textwidth}{0pt}}
    \caption{Arsitektur \textit{Vision Transformer} (ViT): citra dibagi
    menjadi \textit{patch} lalu diproses sebagai urutan token oleh
    \textit{Transformer encoder} \cite{dosovitskiy2021image}}
    \label{fig:vit_architecture}
\end{figure}

\subsubsection{Swin Transformer}
\textit{Swin Transformer} diperkenalkan oleh Liu et al. sebagai
arsitektur \textit{Transformer} hierarkis yang mengintegrasikan konsep
lokalitas dan hierarki dari CNN ke dalam kerangka \textit{Transformer}
\cite{liu2021swin}. Motivasi utama pengembangan ini adalah untuk
mengatasi dua tantangan utama ViT: kompleksitas komputasi yang tumbuh
secara kuadratik terhadap ukuran citra, dan kurangnya representasi
hierarkis yang berguna untuk berbagai tugas penglihatan mesin. \par

Kontribusi utama \textit{Swin Transformer} adalah mekanisme
\textit{shifted window attention} yang membatasi komputasi
\textit{self-attention} hanya dalam jendela lokal yang tidak tumpang
tindih, kemudian menggeser posisi jendela pada lapisan berikutnya
untuk memungkinkan koneksi antar jendela. Pendekatan ini menghasilkan
kompleksitas komputasi yang tumbuh secara linier terhadap ukuran citra,
bukan kuadratik seperti ViT standar. Selain itu, \textit{Swin
Transformer} membangun representasi hierarkis dengan secara bertahap
menggabungkan \textit{patch} yang berdekatan di lapisan yang lebih
dalam, menghasilkan peta fitur dengan resolusi berbeda yang berguna
untuk berbagai tugas \textit{computer vision} \cite{liu2021swin}.
Varian \textit{Swin-Small} yang digunakan dalam penelitian ini memiliki
sekitar 49 juta parameter dengan FLOPs sekitar 8.5 GFLOPs
\cite{liu2021swin}. \par

% TODO: Tambahkan gambar shifted window attention pada Swin Transformer
% Sumber: Liu et al. (2021), Figure 2
\begin{figure}[h]
    \centering
    % \includegraphics[width=0.85\textwidth]{gambar/swin_window.png}
    \fbox{\rule{0pt}{4cm}\rule{0.85\textwidth}{0pt}}
    \caption{Mekanisme \textit{shifted window attention} pada \textit{Swin
    Transformer}: jendela pada lapisan $l$ digeser pada lapisan $l{+}1$
    untuk memungkinkan koneksi antar jendela \cite{liu2021swin}}
    \label{fig:swin_window}
\end{figure}


\subsection{Transfer Learning} \label{II.Teori.Transfer}
\textit{Transfer learning} adalah teknik yang memanfaatkan pengetahuan
yang telah dipelajari dari satu tugas untuk meningkatkan performa pada
tugas yang berbeda namun terkait \cite{pan2010survey}. Dalam konteks
klasifikasi citra bencana, \textit{transfer learning} sangat penting
karena \textit{dataset} anotasi untuk domain spesifik seringkali
terbatas. Model yang telah mempelajari representasi visual umum dari
jutaan citra pada \textit{ImageNet} dapat beradaptasi untuk mengenali
pola-pola spesifik dalam citra bencana dengan hanya beberapa ribu
contoh pelatihan, mengurangi kebutuhan data dan waktu komputasi secara
signifikan. \par

Strategi \textit{fine-tuning} yang digunakan selama \textit{transfer
learning} berpengaruh signifikan terhadap performa akhir model.
\textit{Full fine-tuning} yang langsung melatih seluruh lapisan model
sejak awal berisiko menyebabkan \textit{catastrophic forgetting}
(pelupaan katastrofik, yaitu hilangnya pengetahuan yang telah dipelajari
dari \textit{pre-training} akibat pembaruan bobot yang terlalu agresif)
pada representasi yang sudah dipelajari, terutama ketika \textit{dataset}
target berukuran kecil \cite{kirkpatrick2017catastrophic}. Untuk
mengatasi hal ini, penelitian ini menggunakan strategi
\textit{two-stage fine-tuning} yang membagi proses pelatihan menjadi
dua tahap. Pada Stage 1, \textit{backbone} dibekukan seluruhnya dan
hanya lapisan \textit{classifier head} yang dilatih selama 5 epoch
dengan \textit{learning rate} yang lebih besar, memungkinkan
\textit{head} yang masih acak untuk konvergen terlebih dahulu tanpa
merusak fitur \textit{pretrained}. Pada Stage 2, seluruh lapisan dibuka
dengan \textit{differential learning rate} (laju pembelajaran berbeda
per lapisan) dimana \textit{backbone} menggunakan \textit{learning rate}
10 kali lebih kecil dibandingkan \textit{head}, dengan
\textit{early stopping} (penghentian dini) berbasis \textit{validation
loss} dengan \textit{patience} 5 epoch. Pendekatan ini mencegah
\textit{catastrophic forgetting} sekaligus memungkinkan seluruh model
beradaptasi secara optimal terhadap domain target. \par

% TODO: Buat diagram alur two-stage fine-tuning (dibuat sendiri)
\begin{figure}[h]
    \centering
    % \includegraphics[width=0.85\textwidth]{gambar/two_stage_finetuning.png}
    \fbox{\rule{0pt}{4cm}\rule{0.85\textwidth}{0pt}}
    \caption{Diagram alur \textit{two-stage fine-tuning}: Stage 1
    membekukan \textit{backbone} dan melatih \textit{head}, Stage 2
    membuka seluruh lapisan dengan \textit{differential learning rate}}
    \label{fig:two_stage_finetuning}
\end{figure}


\subsection{Ensemble Learning} \label{II.Teori.Ensemble}
\textit{Ensemble learning} adalah paradigma \textit{machine learning}
yang menggabungkan prediksi dari beberapa model untuk meningkatkan
performa dibandingkan model tunggal manapun \cite{rokach2010ensemble}.
Ide dasarnya adalah bahwa kombinasi dari beberapa model yang saling
melengkapi dapat mengurangi bias dan varians, menghasilkan prediksi
yang lebih \textit{robust} dan akurat. \textit{Heterogeneous ensemble}
yang menggabungkan model dengan arsitektur yang berbeda secara mendasar
cenderung menghasilkan keragaman prediksi yang lebih tinggi, yang
berkontribusi pada peningkatan performa yang lebih signifikan
dibandingkan \textit{homogeneous ensemble} \cite{ganaie2022ensemble}. \par

\subsubsection{Simple Averaging}
\textit{Simple averaging} adalah strategi \textit{ensemble} yang
memberikan bobot sama untuk setiap model. Prediksi akhir dihitung
sebagai rata-rata probabilitas dari seluruh model:

\begin{equationcaptioned}[eq:simple_avg]{
    P_{\text{ensemble}}(c \mid x) = \frac{1}{M} \sum_{m=1}^{M}
    P_m(c \mid x)
}{
    Simple Averaging Ensemble
}
\end{equationcaptioned}

dimana $M$ adalah jumlah model, $P_m(c \mid x)$ adalah probabilitas
prediksi model ke-$m$ untuk kelas $c$ pada masukan $x$, dan kelas
akhir dipilih berdasarkan probabilitas tertinggi. \par

\subsubsection{Weighted Voting}
\textit{Weighted voting} memberikan bobot berbeda kepada masing-masing
model berdasarkan performa validasinya, sehingga model yang lebih baik
mendapat pengaruh lebih besar dalam prediksi akhir:

\begin{equationcaptioned}[eq:weighted_vote]{
    P_{\text{ensemble}}(c \mid x) = \frac{\sum_{m=1}^{M} w_m \cdot
    P_m(c \mid x)}{\sum_{m=1}^{M} w_m}
}{
    Weighted Voting Ensemble
}
\end{equationcaptioned}

dimana $w_m$ adalah bobot model ke-$m$ yang ditetapkan berdasarkan
akurasi validasinya. \par

\subsubsection{Stacking dengan Logistic Regression}
\textit{Stacking} adalah strategi \textit{ensemble} yang melatih sebuah
\textit{meta-learner} untuk mempelajari cara terbaik menggabungkan
prediksi dari model-model dasar \cite{ganaie2022ensemble}. Berbeda
dengan \textit{simple averaging} dan \textit{weighted voting} yang
menggunakan aturan kombinasi yang telah ditetapkan, \textit{stacking}
mempelajari kombinasi tersebut secara adaptif dari data. \par

Dalam penelitian ini, \textit{meta-learner} yang digunakan adalah
\textit{Logistic Regression} yang dilatih pada vektor probabilitas yang
dihasilkan oleh keempat model dasar pada \textit{validation set}. Vektor
fitur untuk setiap sampel dibentuk dengan menggabungkan (\textit{concatenate})
vektor probabilitas dari seluruh model dasar, sehingga untuk $M$ model
dan $C$ kelas, dimensi fitur masukan \textit{meta-learner} adalah
$M \times C$. \textit{Meta-learner} kemudian dievaluasi pada
\textit{test set} menggunakan probabilitas dari model-model dasar yang
telah dilatih. Penggunaan \textit{Logistic Regression} sebagai \textit{meta-learner}
dipilih karena sifatnya yang mudah diinterpretasi dan terbukti efektif
untuk masalah kombinasi prediksi probabilistik
\cite{ganaie2022ensemble}. \par


\subsection{Augmentasi Data} \label{II.Teori.Augmentasi}
Augmentasi data adalah teknik untuk meningkatkan keragaman data
pelatihan dengan membuat variasi artifisial dari citra yang ada tanpa
mengubah labelnya \cite{shorten2019survey}. Teknik ini bertujuan
mengurangi risiko \textit{overfitting} (kondisi di mana model terlalu
menyesuaikan diri dengan data pelatihan sehingga gagal menggeneralisasi
pada data baru) dan meningkatkan kemampuan generalisasi model. Teknik augmentasi yang digunakan dalam penelitian
ini meliputi \textit{random resized crop} untuk variasi skala dan posisi,
\textit{horizontal flip} untuk variasi orientasi, rotasi acak,
\textit{color jitter} untuk variasi kecerahan dan kontras, serta
\textit{random erasing} yang secara acak menutupi sebagian wilayah
citra untuk memaksa model belajar fitur yang lebih \textit{robust}.
Teknik-teknik ini membantu model menjadi lebih \textit{robust} terhadap
variasi yang umum terjadi pada citra media sosial. \par


\subsection{Optimizer dan Learning Rate Schedule}
\label{II.Teori.Optimizer}

\subsubsection{AdamW}
AdamW adalah varian dari algoritma \textit{optimizer} Adam yang
memisahkan \textit{weight decay} (peluruhan bobot, yaitu penalti
regularisasi yang mendorong bobot model tetap kecil) dari pembaruan
gradien \cite{loshchilov2019adamw}. Berbeda dengan Adam standar yang
menggabungkan \textit{weight decay} ke dalam gradien sebelum normalisasi
adaptif, AdamW menerapkan \textit{weight decay} langsung pada parameter
model setelah pembaruan adaptif. Pemisahan ini menghasilkan regularisasi
yang lebih tepat dan terbukti lebih efektif untuk melatih model berbasis
\textit{Transformer} seperti ViT dan \textit{Swin Transformer}. \par

\subsubsection{Cosine Annealing}
\textit{Cosine annealing} adalah strategi penjadwalan \textit{learning
rate} yang menurunkan \textit{learning rate} mengikuti fungsi kosinus
dari nilai awal menuju nilai minimum \cite{loshchilov2017sgdr}.
Penjadwalan ini memungkinkan model menjelajahi ruang parameter secara
luas pada awal pelatihan dan kemudian melakukan konvergensi yang lebih
halus menuju solusi optimal, menghasilkan proses pelatihan yang lebih
stabil dan performa akhir yang lebih baik. \par


\subsection{Loss Function} \label{II.Teori.Loss}

\subsubsection{Cross-Entropy Loss}
\textit{Cross-entropy loss} adalah fungsi kerugian yang paling umum
digunakan untuk tugas klasifikasi \cite{goodfellow2016deep}. Fungsi
ini mengukur perbedaan antara distribusi probabilitas prediksi model
dengan distribusi label sebenarnya:

\begin{equationcaptioned}[eq:cross_entropy]{
    \mathcal{L}_{CE} = -\frac{1}{N} \sum_{i=1}^{N} \sum_{c=1}^{C}
    y_{i,c} \log(\hat{p}_{i,c})
}{
    Cross-Entropy Loss
}
\end{equationcaptioned}

dimana $N$ adalah jumlah sampel, $C$ adalah jumlah kelas, $y_{i,c}$
adalah label \textit{one-hot} yang bernilai 1 jika sampel ke-$i$
termasuk kelas $c$ dan 0 sebaliknya, serta $\hat{p}_{i,c}$ adalah
probabilitas prediksi model untuk sampel ke-$i$ pada kelas $c$.
Penelitian ini menggunakan \textit{label smoothing} (pelunakan label)
dengan nilai 0.1 yang mendistribusikan sebagian kecil probabilitas
kepada kelas yang salah, berfungsi sebagai regularisasi untuk mencegah
model terlalu percaya diri terhadap prediksinya. \par

\subsubsection{Weighted Cross-Entropy Loss}
Pada \textit{dataset} dengan distribusi kelas yang tidak seimbang,
\textit{weighted cross-entropy loss} memberikan bobot yang lebih besar
pada kelas minoritas untuk mencegah model bias terhadap kelas mayoritas:

\begin{equationcaptioned}[eq:weighted_cross_entropy]{
    \mathcal{L}_{WCE} = -\frac{1}{N} \sum_{i=1}^{N} \sum_{c=1}^{C}
    w_c \cdot y_{i,c} \log(\hat{p}_{i,c})
}{
    Weighted Cross-Entropy Loss
}
\end{equationcaptioned}

dimana $w_c$ adalah bobot untuk kelas $c$ yang dihitung sebagai
kebalikan dari frekuensi kemunculan kelas tersebut dalam data
pelatihan. \par


\subsection{Metrik Evaluasi} \label{II.Teori.Metrik}

\subsubsection{Accuracy}
\textit{Accuracy} mengukur proporsi prediksi yang benar dari seluruh
prediksi yang dibuat oleh model:

\begin{equationcaptioned}[eq:accuracy]{
    \text{Accuracy} = \frac{TP + TN}{TP + TN + FP + FN}
}{
    Accuracy
}
\end{equationcaptioned}

dimana $TP$ (\textit{True Positive}), $TN$ (\textit{True Negative}),
$FP$ (\textit{False Positive}), dan $FN$ (\textit{False Negative})
masing-masing adalah jumlah prediksi positif benar, prediksi negatif
benar, prediksi positif salah, dan prediksi negatif salah. Meskipun
intuitif, \textit{accuracy} dapat menyesatkan pada \textit{dataset}
yang tidak seimbang. \par

\subsubsection{Precision dan Recall}
\textit{Precision} mengukur proporsi prediksi positif yang benar dari
seluruh prediksi positif yang dibuat oleh model:

\begin{equationcaptioned}[eq:precision]{
    \text{Precision} = \frac{TP}{TP + FP}
}{
    Precision
}
\end{equationcaptioned}

\textit{Recall} mengukur proporsi \textit{instance} positif yang benar
diidentifikasi dari seluruh \textit{instance} positif sebenarnya:

\begin{equationcaptioned}[eq:recall]{
    \text{Recall} = \frac{TP}{TP + FN}
}{
    Recall
}
\end{equationcaptioned}

\textit{Recall} sangat penting dalam deteksi bencana dimana kehilangan
informasi kritis dapat berakibat fatal. \par

\subsubsection{F1-Score}
\textit{F1-score} adalah \textit{harmonic mean} dari \textit{precision}
dan \textit{recall} \cite{powers2011evaluation}:

\begin{equationcaptioned}[eq:f1score]{
    F1 = 2 \times \frac{\text{Precision} \times \text{Recall}}{
    \text{Precision} + \text{Recall}}
}{
    F1-Score
}
\end{equationcaptioned}

Dalam konteks klasifikasi multi-kelas, \textit{macro-averaging}
memberikan bobot sama untuk setiap kelas tanpa mempertimbangkan
jumlah sampel per kelas:

\begin{equationcaptioned}[eq:f1macro]{
    F1_{\text{macro}} = \frac{1}{C} \sum_{c=1}^{C} F1_c
}{
    F1-Score Macro-Averaging
}
\end{equationcaptioned}

\textit{Weighted-averaging} memberikan bobot sesuai jumlah
\textit{instance} per kelas:

\begin{equationcaptioned}[eq:f1weighted]{
    F1_{\text{weighted}} = \frac{1}{N} \sum_{c=1}^{C} n_c \cdot F1_c
}{
    F1-Score Weighted-Averaging
}
\end{equationcaptioned}

dimana $n_c$ adalah jumlah \textit{instance} pada kelas $c$ dan
$N = \sum_{c=1}^{C} n_c$. \textit{Macro F1} digunakan sebagai metrik
utama dalam penelitian ini karena memberikan bobot yang sama untuk
setiap kelas, sehingga mencerminkan performa model pada kelas minoritas
secara adil dan memungkinkan perbandingan yang lebih bermakna pada
\textit{dataset} yang tidak seimbang. \par

\subsubsection{AUC-ROC}
\textit{Area Under the Receiver Operating Characteristic Curve}
(AUC-ROC) mengukur kemampuan model untuk membedakan antar kelas secara
keseluruhan, terlepas dari nilai ambang batas keputusan yang digunakan
\cite{powers2011evaluation}. Nilai AUC-ROC berkisar antara 0 hingga 1,
dimana 1 menunjukkan pemisahan kelas yang sempurna dan 0.5 setara
dengan tebakan acak. Untuk klasifikasi biner, AUC-ROC dihitung
langsung dari probabilitas kelas positif. Untuk klasifikasi multi-kelas,
penelitian ini menggunakan pendekatan \textit{One-vs-Rest} (OvR) dengan
\textit{macro-averaging}, dimana AUC-ROC dihitung untuk setiap kelas
terhadap semua kelas lainnya kemudian dirata-ratakan:

\begin{equationcaptioned}[eq:aucroc]{
    \text{AUC-ROC}_{\text{macro}} = \frac{1}{C} \sum_{c=1}^{C}
    \text{AUC}_c
}{
    AUC-ROC Macro OvR
}
\end{equationcaptioned}

\subsubsection{Confusion Matrix}
\textit{Confusion matrix} adalah tabel berukuran $C \times C$ yang
mendeskripsikan performa model klasifikasi secara mendetail
\cite{powers2011evaluation}. Elemen pada baris $i$ dan kolom $j$
menunjukkan jumlah \textit{instance} dari kelas $i$ yang diprediksi
sebagai kelas $j$. Matriks ini memberikan visualisasi yang jelas tentang
jenis kesalahan yang dibuat model dan membantu identifikasi kelas mana
yang sulit diklasifikasikan. \par

\subsubsection{Metrik Efisiensi Komputasi}
Evaluasi efisiensi komputasi dilakukan sebagai analisis pendukung pada
Task 1 untuk memberikan gambaran biaya komputasi dari masing-masing
konfigurasi model \cite{menghani2023efficient}. \textit{Inference time}
adalah rata-rata waktu yang dibutuhkan model untuk memproses satu citra
dalam milidetik (ms):

\begin{equationcaptioned}[eq:inference_time]{
    t_{\text{avg}} = \frac{1}{B} \sum_{i=1}^{B} t_i
}{
    Rata-rata Inference Time
}
\end{equationcaptioned}

dimana $B$ adalah jumlah citra dalam \textit{batch} dan $t_i$ adalah
waktu inferensi untuk citra ke-$i$. \textit{Memory usage} adalah
konsumsi memori GPU maksimum selama inferensi dalam megabyte (MB).
FLOPs adalah jumlah operasi matematika \textit{floating-point} yang
dibutuhkan untuk memproses satu citra dalam GigaFLOPs (GFLOPs),
memberikan indikasi kompleksitas komputasi yang independen dari
perangkat keras spesifik:

\begin{equationcaptioned}[eq:flops]{
    \text{FLOPs} = \sum_{\text{layers}} \text{FLOPs}_{\text{layer}}
}{
    Total FLOPs
}
\end{equationcaptioned}

Untuk \textit{convolutional layer}, FLOPs dihitung sebagai:

\begin{equationcaptioned}[eq:flops_conv]{
    \text{FLOPs}_{\text{conv}} = 2 \times H_{\text{out}} \times
    W_{\text{out}} \times C_{\text{in}} \times C_{\text{out}} \times K^2
}{
    FLOPs Convolutional Layer
}
\end{equationcaptioned}

dimana $H_{\text{out}}$ dan $W_{\text{out}}$ adalah dimensi output,
$C_{\text{in}}$ dan $C_{\text{out}}$ adalah jumlah \textit{channel}
input dan output, dan $K$ adalah ukuran kernel. \par


\subsection{Dataset CrisisMMD} \label{II.Teori.Dataset}
\textit{CrisisMMD} adalah \textit{dataset} multimodal yang dikurasi
khusus untuk penelitian klasifikasi informasi media sosial dalam konteks
respons bencana \cite{alam2018crisismmd}. \textit{Dataset} ini
mengandung lebih dari 18.000 citra yang dikumpulkan dari tujuh kejadian
bencana alam yang semuanya terjadi pada tahun 2017, yaitu
\textit{Hurricane Harvey}, \textit{Hurricane Irma},
\textit{Hurricane Maria}, gempa bumi Mexico, kebakaran hutan California,
gempa bumi Iraq-Iran, dan banjir Sri Lanka. Setiap \textit{instance}
dalam \textit{dataset} telah dianotasi secara manual oleh beberapa
anotator untuk memastikan kualitas label. \par

\textit{Dataset} menyediakan tiga tugas klasifikasi utama yang menjadi
fokus penelitian ini: \par

\textbf{Task 1 -- \textit{Informative Classification}:} Klasifikasi
citra sebagai \textit{Informative} (mengandung informasi berguna untuk
respons bencana) atau \textit{Not Informative} (tidak relevan atau
bukan citra bencana). Ini adalah tugas klasifikasi biner. \par

\textbf{Task 2 -- \textit{Humanitarian Categories}:} Klasifikasi citra
ke dalam delapan kategori: tidak bersifat kemanusiaan
(\textit{not humanitarian}), kerusakan infrastruktur dan utilitas,
informasi relevan lainnya, upaya penyelamatan atau donasi, individu
terdampak, kerusakan kendaraan, orang terluka atau meninggal, serta
orang hilang atau ditemukan. Ini adalah tugas klasifikasi multi-kelas
dengan distribusi kelas yang tidak seimbang. \par

\textbf{Task 3 -- \textit{Damage Severity Assessment}:} Penilaian
tingkat kerusakan yang ditampilkan dalam gambar ke dalam tiga kategori:
sedikit atau tidak ada kerusakan (\textit{little or no damage}),
kerusakan ringan (\textit{mild damage}), dan kerusakan parah
(\textit{severe damage}). Task 3 merupakan tugas berbasis citra saja
dengan total sampel yang lebih kecil dibandingkan Task 1 dan Task 2. \par

Ketiga tugas menggunakan \textit{split} resmi yang disediakan oleh
pembuat \textit{dataset} dengan rasio 70\% \textit{training set},
15\% \textit{validation set}, dan 15\% \textit{test set}.
Pembagian ini dilakukan secara \textit{stratified} sehingga proporsi
kelas pada setiap \textit{split} konsisten dengan distribusi
keseluruhan. Rincian pembagian data disajikan pada
Tabel~\ref{table:crisismmd_task1},
Tabel~\ref{table:crisismmd_task2}, dan
Tabel~\ref{table:crisismmd_task3}. \par

\begin{table}[h]
    \centering
    \caption{Pembagian data \textit{CrisisMMD} untuk Task 1
    \cite{alam2018crisismmd}}
    \label{table:crisismmd_task1}
    \begin{tabular}{|l|r|r|r|r|}
        \hline
        \textbf{Kelas} &
        \textbf{Train} &
        \textbf{Dev} &
        \textbf{Test} &
        \textbf{Total} \\
        \hline
        Informative     & 7.059 & 1.164 & 1.151 & 9.374 \\
        Not Informative & 6.549 & 1.073 & 1.086 & 8.708 \\
        \hline
        \textbf{Total}  & \textbf{13.608} & \textbf{2.237} &
                          \textbf{2.237}  & \textbf{18.082} \\
        \hline
    \end{tabular}
\end{table}

\begin{table}[h]
    \centering
    \caption{Pembagian data \textit{CrisisMMD} untuk Task 2
    \cite{alam2018crisismmd}}
    \label{table:crisismmd_task2}
    \begin{tabular}{|l|r|r|r|r|}
        \hline
        \textbf{Kelas} &
        \textbf{Train} &
        \textbf{Dev} &
        \textbf{Test} &
        \textbf{Total} \\
        \hline
        Not humanitarian                       & 6.565 & 1.088 & 1.055 & 8.708 \\
        Infrastructure and utility damage      & 2.761 &   405 &   458 & 3.624 \\
        Other relevant information             & 1.860 &   338 &   331 & 2.529 \\
        Rescue, volunteering or donation       & 1.694 &   281 &   256 & 2.231 \\
        Affected individuals                   &   397 &    77 &    88 &   562 \\
        Vehicle damage                         &   230 &    38 &    36 &   304 \\
        Injured or dead people                 &    92 &     8 &    10 &   110 \\
        Missing or found people                &     9 &     2 &     3 &    14 \\
        \hline
        \textbf{Total}                         & \textbf{13.608} & \textbf{2.237} &
                                                 \textbf{2.237}  & \textbf{18.082} \\
        \hline
    \end{tabular}
\end{table}

\begin{table}[h]
    \centering
    \caption{Pembagian data \textit{CrisisMMD} untuk Task 3
    \cite{alam2018crisismmd}}
    \label{table:crisismmd_task3}
    \begin{tabular}{|l|r|r|r|r|}
        \hline
        \textbf{Kelas} &
        \textbf{Train} &
        \textbf{Dev} &
        \textbf{Test} &
        \textbf{Total} \\
        \hline
        Severe damage        & 1.548 & 332 & 332 & 2.212 \\
        Mild damage          &   587 & 126 & 126 &   839 \\
        Little or no damage  &   333 &  71 &  71 &   475 \\
        \hline
        \textbf{Total}       & \textbf{2.468} & \textbf{529} &
                               \textbf{529}   & \textbf{3.526} \\
        \hline
    \end{tabular}
\end{table}
